\manualchapter{Patching and operation}{sec:operation}
\subsection{Pulse duet}
Use the first two pulse outputs separately, tune the second slightly sharp
with MOD and send both through an external mixer. A pitch cable in the first
V/OCT input controls both until the second input is patched. Different PW
settings make the two voices easier to distinguish.
\subsection{Noise percussion and sync}
Patch a short envelope to the noise AMP input and monitor NOISE. Patch the
other outputs, or lower their levels, when you want to exclude their normalled
mix. Change Period and LFSR for different textures. A trigger into noise SYNC
makes repeated hits start from a consistent noise phase.
\subsection{Polyphony and saved patches}
The widest input selects 1--16 channels. This module reuses a mono CV on every
active channel; polyphonic CV controls channels independently. Pitch, FM and
sync are processed each sample; volume, duty and noise controls update every
16 samples. Changing sample rate adjusts clock quantization and can shift
extreme register-limited pitches slightly.

Rack saves controls and cables. There is no custom oscillator-state save or
extra module context menu. Initialize restores controls and resets the chip;
reloading a patch does not resume oscillator phase at its saved point.
\subsection{Troubleshooting}
If several sounds appear at one jack, check output normalling before changing
oscillator settings. If an envelope sounds stepped, lower its speed or use an
external VCA. If the triangle does not fade smoothly, use an external VCA
rather than treating its chip level as a linear amplitude control. If a gate
does not restart a pulse, use an amplitude envelope; only Triangle and Noise
have SYNC inputs.
